
\subsection*{6.3.}
Montrons d'abord comment 6.2.3 découle de 6.2.2. Pour cela il suffit évidemment de démontrer (6.2.5) pour les
@@M0@@ @@M1@@, donc on peut supposer @@M2@@, @@M3@@. D'après 6.2.2 @@M4@@ se factorise par @@M5@@ i.e. il existe un diagramme commutatif
@@M6@@
où @@M7@@ est l'épimorphisme canonique et @@M8@@ une représentation de @@M9@@ dans @@M10@@. Mais l'homomorphisme
@@M11@@
possède une section, de sorte que le monomorphisme
@@M12@@
induit par det est un isomorphisme. D'autre part on sait (SGA 3 VIII 1.3) que toute représentation de @@M13@@ dans lui-même est de la forme
@@M14@@
où @@M15@@ est une section de @@M16@@, d'où l'assertion.

Le corollaire 6.2.6 est une simple reformulation de 6.2.3. Notons que 6.2.3 découle aussi immédiatement de SGA 3 XXII 6.2.1, mais nous voulons lui donner ici une démonstration élémentaire via 6.2.2.

